\section*{第 \ref{ch:b1:photosynthesis} 章——光合作用与自养}

\begin{solution}{exo:b1:photosynthesis:1}
外膜和内膜，以及
内部的\omterm{def:b1:photosynthesis:chloroplast}{类囊体}膜。\omterm{prop:b1:photosynthesis:lightreactions}{光反应}：在\omterm{def:b1:photosynthesis:chloroplast}{类囊体}膜上进行（质子被泵入
\omterm{def:b1:photosynthesis:chloroplast}{类囊体}腔中）。\omterm{def:b1:photosynthesis:fixation}{卡尔文循环}：在\omterm{def:b1:photosynthesis:chloroplast}{基质}中进行。
\end{solution}

\begin{solution}{exo:b1:photosynthesis:2}
$2\,\mathrm{H_2O} + 2\,\mathrm{NADP^+} + 3\,\mathrm{ADP} + 3\,\mathrm{P_i}
\to \mathrm{O_2} + 2\,\mathrm{NADPH} + 2\,\mathrm{H^+} + 3\,\mathrm{ATP}$
需八个光子。氧气来自水：以
$\mathrm{H_2^{18}O}$ 喂养的藻类释放 $\mathrm{^{18}O_2}$；以
$\mathrm{C^{18}O_2}$ 喂养的则不释放（鲁本和卡门）。
\end{solution}

\begin{solution}{exo:b1:photosynthesis:3}
\omterm{def:b1:photosynthesis:fixation}{固定}（\omterm{def:b1:photosynthesis:fixation}{RuBisCO}）、还原（生成 G3P）、再生（RuBP）。每
$\mathrm{CO_2}$ 需 3 个 \omterm{def:b1:water-small-molecules:atp}{ATP} 和 2 个 NADPH。
\end{solution}

\begin{solution}{exo:b1:photosynthesis:4}
在 \qty{430}{nm}（蓝光）和 \qty{660}{nm}（红光）附近吸收最强。在 \qty{500}{nm}
左右，\omterm{def:b1:photosynthesis:pigments}{类胡萝卜素}吸收光能并将其传给\omterm{def:b1:photosynthesis:pigments}{叶绿素}，因此
即使\omterm{def:b1:photosynthesis:pigments}{叶绿素}自身在此吸收很少，仍释放出氧气。
\end{solution}

\begin{solution}{exo:b1:photosynthesis:5}
远红光（\qty{700}{nm}）单独照射时，几乎不驱动光合作用；
与 \qty{680}{nm} 光同时照射时，其速率超过两者之和。
两个吸收峰不同的色素系统必须串联协作，
各需其自身的光子：\omterm{def:b1:photosynthesis:zscheme}{光系统} I（P700）和
\omterm{def:b1:photosynthesis:zscheme}{光系统} II（P680）。
\end{solution}

\begin{solution}{exo:b1:photosynthesis:6}
每摩尔质子 $0.18\times 96\,500 = \qty{17.4}{kJ/mol}$；至少需 $50/17.4 = 2.9$
个质子；合成酶用了四个，即
效率为 \qty{72}{\%}。
\end{solution}

\begin{solution}{exo:b1:photosynthesis:7}
这证明，跨\omterm{def:b1:photosynthesis:chloroplast}{类囊体}膜的质子梯度
足以驱动 \omterm{def:b1:water-small-molecules:atp}{ATP} 合成，无需光照，也无需
电子传递：二者通过梯度偶联。但它本身
并不能证明叶中\omterm{prop:b1:photosynthesis:lightreactions}{光反应}以这种方式合成 \omterm{def:b1:water-small-molecules:atp}{ATP}（此需
解偶联剂及 \omterm{def:b1:water-small-molecules:ph}{pH} 测量）。加入解偶联剂后，梯度
在合成酶利用之前即已消散，无 \omterm{def:b1:water-small-molecules:atp}{ATP} 生成。
\end{solution}

\begin{solution}{exo:b1:photosynthesis:8}
$3\times 5 + 3\times 1 = 18$ 个碳；六个 3-PGA $= 18$；六个 G3P $= 18$；
移出一个 G3P（3），五个 G3P（15）再生三个 RuBP（15）：$3 +
15 = 18$。平衡。
\end{solution}

\begin{solution}{exo:b1:photosynthesis:9}
无 $\mathrm{CO_2}$ 时 \omterm{def:b1:photosynthesis:fixation}{RuBisCO} 停止工作：RuBP 不再被消耗
（上升），3-PGA 不再生成（下降），而光照仍在
将其还原掉。无光时没有 \omterm{def:b1:water-small-molecules:atp}{ATP} 和 NADPH：3-PGA 不再被还原
（上升），RuBP 不再得到再生（下降），而\omterm{def:b1:photosynthesis:fixation}{固定}
仍短暂进行。
\end{solution}

\begin{solution}{exo:b1:photosynthesis:10}
$E = hc/\lambda = 6.63\times 10^{-34}\times 3\times 10^8/6.8\times
10^{-7} = \qty{2.92e-19}{J}$；每摩尔 \qty{176}{kJ}。48 个光子：
\qty{8450}{kJ}；效率 $2870/8450 = \qty{34}{\%}$。作物损失的
未被吸收的光（一半光谱、反射、透射及
植株间的空隙）、光呼吸、叶片和茎的呼吸作用、
根的呼吸作用、全日照下酶的光饱和，以及
无叶的季节。
\end{solution}

\begin{solution}{exo:b1:photosynthesis:11}
每 100 次羧化中，有 33 次加氧，释放 $16.7$ 个碳：净
$83.3$，损失 \qty{17}{\%}。避免此损失的代价为：2 个 \omterm{def:b1:water-small-molecules:atp}{ATP} $\approx
0.2$ 个碳（每 $\mathrm{CO_2}$），即 \qty{20}{\%}；C4 设计只在
损失超过约五分之一时才划算——即炎热干燥的条件下，
而非凉爽的条件下。
\end{solution}

\begin{solution}{exo:b1:photosynthesis:12}
两者都利用膜结构，借电子传递的能量泵送质子，
并由 \omterm{def:b1:water-small-molecules:atp}{ATP} 合成酶消耗该梯度：机制
相同，两种合成酶为同源。但电子的流向
相反——在\omterm{def:b1:photosynthesis:chloroplast}{叶绿体}中由水流向 NADPH，靠光能提升，
在\omterm{def:b1:eukaryotic-cell:mitochondrion}{线粒体}中由 NADH 流向氧，顺势下降
（见\cref{ch:b1:respiration-fermentation}）——且两者的碳循环
并非互为逆过程：\omterm{def:b1:photosynthesis:fixation}{卡尔文循环}与克雷布斯循环不共享
任何酶，$\mathrm{CO_2}$ 的还原消耗 NADPH，而其
释放则使用 $\mathrm{NAD^+}$。总反应式互为逆向；
机制是共同遗产，被用于两个方向，
在碳端的化学过程则各不相同。
\end{solution}

\begin{solution}{pb:b1:photosynthesis:1}
\textbf{1.} $hc/\lambda = \qty{2.92e-19}{J}$；$\times N_A = \qty{176}{kJ/mol}$。
\textbf{2.} \qty{700}{nm}：\qty{171}{kJ/mol}；\qty{450}{nm}：
\qty{266}{kJ/mol}。蓝光光子将\omterm{def:b1:photosynthesis:pigments}{叶绿素}激发至较高能态，
该态在皮秒内衰变至红光光子所到达的同一最低激发态，
多余能量以热的形式散失，化学过程中两者
都只相当于一个光子。
\textbf{3.} 每焦耳：$1/2.92\times 10^{-19} = \num{3.4e18}$ 个光子。
\textbf{4.} $\Delta E = 0.82 - (-0.32) = \qty{1.14}{V}$：$\Delta G =
96\,500\times 1.14 = \qty{110}{kJ}$，每摩尔电子。
\textbf{5.} 两个光子：约 $176 + 171 = \qty{347}{kJ}$；储存
\qty{110}{kJ}，即 \qty{32}{\%}。
\textbf{6.} 每个 $\mathrm{O_2}$ 对应四个电子（两个水分子），
生成两个 NADPH：八个光子。
\textbf{7.} 12 个 NADPH：24 个电子，48 个光子。
\textbf{8.} $2.303\times 8.314\times 298\times 3/96\,500 = \qty{0.177}{V}$；
每摩尔质子 $\qty{17.1}{kJ}$。
\textbf{9.} 共 $4 + 8 = 12$ 个质子（每 $\mathrm{O_2}$）：合成 3 个 \omterm{def:b1:water-small-molecules:atp}{ATP}。
\textbf{10.} $6\times 3 = 18$ 个 \omterm{def:b1:water-small-molecules:atp}{ATP}：恰好等于该核算方式下循环的需求，
但实际每 $\mathrm{O_2}$ 的产量接近 2.6
（\omterm{def:b1:photosynthesis:chloroplast}{叶绿体}中每个 \omterm{def:b1:water-small-molecules:atp}{ATP} 耗 \qty{4.7}{H^+}），其余部分
由绕\omterm{def:b1:photosynthesis:zscheme}{光系统} I 的环式电子流补足。
\textbf{11.} $4\times 17.1 = \qty{68}{kJ}$，对应一个 \qty{50}{kJ} 的 \omterm{def:b1:water-small-molecules:atp}{ATP}：
\qty{73}{\%}。
\textbf{12.} $10^{-5}\,\mathrm{mol/L}\times 10^{-15}\,\mathrm{L}\times
6\times 10^{23} = 6$ 个游离质子存于腔中，而每秒经合成酶通过的 $10^5$
个质子：通量由结合在
腔内\omterm{def:b1:proteins:peptide}{蛋白质}与脂质缓冲基团上的质子承载，这些基团
随质子离开而迅速将其释放。
\textbf{13.} $18\times 50 + 12\times 220 = 900 + 2640 = \qty{3540}{kJ}$
对应储存的 \qty{2870}{kJ}：\qty{670}{kJ}（约五分之一）以热的形式耗散于
循环的反应中，这些反应须为放能反应方能进行。
\textbf{14.} $48\times 176 = \qty{8450}{kJ}$；$2870/8450 = \qty{34}{\%}$。
\textbf{15.} $0.34\times 0.45\times 0.9 = \qty{13.8}{\%}$。
\textbf{16.} $400\times 0.005\times 36\,000 = \qty{72}{kJ}$；葡萄糖
$0.138\times 72\,000/2\,870\,000 = \qty{3.5e-3}{mol} = \qty{0.62}{g}$。
\textbf{17.} 净得葡萄糖 $0.6\times 0.62 = \qty{0.37}{g}$，
另得碳 $\qty{0.15}{g}$（$72/180$）。
\textbf{18.} 光照落在植株之间或土壤上；叶片处于
其他叶片的遮荫之下，以低于容量的水平工作；全日照下的叶片
已饱和（酶无法利用全部光子）；根、茎的呼吸作用
及夜间呼吸；光呼吸；冠层闭合前的数周
以及冠层衰老之后的数周。
\textbf{19.} 33 次加氧，损失 16.7 个碳，净得 83.3 个碳。
\textbf{20.} 回收 $33\times 3.5 = 117$ 个 \omterm{def:b1:water-small-molecules:atp}{ATP}；循环 300 个 \omterm{def:b1:water-small-molecules:atp}{ATP}；合计
417 个 \omterm{def:b1:water-small-molecules:atp}{ATP}，对应 83.3 个碳：每个净碳 5.0 个 \omterm{def:b1:water-small-molecules:atp}{ATP}。
\textbf{21.} $3 + 2 = 5$ 个 \omterm{def:b1:water-small-molecules:atp}{ATP} 每个碳，无碳损失——两者
在 \omterm{def:b1:water-small-molecules:atp}{ATP} 上相等，但 C3 植物还损失了 \qty{17}{\%} 的羧化能力；
在 \qty{25}{\celsius} 下大致持平。
\textbf{22.} 50 次加氧，损失 25 个碳，净得 75 个碳；回收 175 个 \omterm{def:b1:water-small-molecules:atp}{ATP}；
每个净碳 $475/75 = 6.3$ 个 \omterm{def:b1:water-small-molecules:atp}{ATP}，对比 C4 植物的 5 个，但
C3 植物损失了四分之一的碳：C4 显著胜出。
\textbf{23.} 将 $\mathrm{CO_2}$ 增至三倍，使 \omterm{def:b1:photosynthesis:fixation}{RuBisCO} 的竞争
转向羧化，减少光呼吸，并提高 C3 植物的净\omterm{def:b1:photosynthesis:fixation}{固定}；
C4 植物已借浓缩的
$\mathrm{CO_2}$ 使 \omterm{def:b1:photosynthesis:fixation}{RuBisCO} 饱和，除少量节水外别无所得。
\textbf{24.} \qty{1}{g} 碳即 \qty{0.083}{mol}：C3 植物需水 $250\times 0.083 = \qty{21}{mol}$，即 \qty{375}{g}；CAM 植物为 \qty{37}{g}。
\textbf{25.} 每个 $\mathrm{CO_2}$ 约 8 个光子（每葡萄糖 48 个）；
从被吸收光子到葡萄糖的效率为 \qty{34}{\%}；从入射阳光到总糖约
\qty{14}{\%}，到叶片净糖约 \qty{8}{\%}——这一数字
经整株植物整季的损耗降至
\qty{1}{\%}。
\end{solution}
